\chapter{Impulse and Frequency Response}

In the last two chapters we learned that a linear, time-invariant (LTI)
system is completely described by a single sequence: its response to a unit
impulse. We called that sequence the \emph{impulse response} $h[n]$, and we
saw that the output of the system is the convolution of the input with $h$.
That is a central fact, but on its own it can feel a little
opaque. If I hand you a list of forty numbers and tell you ``this is your
filter,'' can you tell me what it \emph{does} to a sound? To a tremor in
your wrist? To a 60~Hz hum on a sensor wire?

The answer is yes, and the key is to look at the same system through a second
window. The impulse response lives in \emph{time}: it tells you how the system
spreads a single poke out over the following samples. The \emph{frequency
response} lives in \emph{frequency}: it tells you, for each pure tone you could
feed in, how much that tone is amplified or attenuated and how much it is
delayed. These two descriptions are the same object seen
from two sides, joined by the Fourier transform. Learning to flip between them
fluently --- ``what does this $h[n]$ do to a 100~Hz sine?'' and ``what $h[n]$
gives me \emph{that} curve?'' --- is most of what it means to think like a
filter designer.

This chapter builds the frequency response geometrically: a filter is a machine
that takes a complex exponential in and gives the \emph{same} complex
exponential back out, only rescaled by a complex number. That single number,
as the input frequency sweeps, traces out the complete frequency response of the system.

\section{The defining experiment: complex exponentials as eigenfunctions}
\label{ch14:eigen}

Let us run the simplest possible test. Feed an LTI system a single pure complex
exponential that has existed for all time,
\[
  x[n] = e^{\jj \omega n}, \qquad n \in \Z,
\]
where $\omega$ is the angular frequency in radians per sample. Because the
system is LTI, its output is the convolution of $x$ with the impulse response
$h[n]$:
\[
  y[n] = (h \conv x)[n]
       = \sum_{k=-\infty}^{\infty} h[k]\, x[n-k]
       = \sum_{k=-\infty}^{\infty} h[k]\, e^{\jj \omega (n-k)}.
\]
Now pull the part that does not depend on $k$ out of the sum:
\[
  y[n] = \underbrace{e^{\jj \omega n}}_{\text{the input}}
         \cdot \underbrace{\sum_{k=-\infty}^{\infty} h[k]\, e^{-\jj \omega k}}_{\text{a number}} .
\]
The output is the \emph{same} complex exponential
$e^{\jj\omega n}$ we put in, multiplied by a constant that does not depend on
$n$ at all. The constant is built only from $h$ and the chosen frequency
$\omega$. We give it a name:
\begin{equation}
  H(e^{\jj\omega}) \;=\; \sum_{k=-\infty}^{\infty} h[k]\, e^{-\jj \omega k}.
  \label{ch14:Hdef}
\end{equation}
So the experiment yields (Figure~\ref{ch14:fig-eigen})
\begin{equation}
  x[n] = e^{\jj\omega n} \;\;\longrightarrow\;\; y[n] = H(e^{\jj\omega})\, e^{\jj\omega n}.
  \label{ch14:eigenrelation}
\end{equation}

\begin{figure}[H]
\centering
\begin{tikzpicture}[node distance=8mm, >=Latex]
\node (in) {$e^{\jj\omega n}$};
\node[draw, rounded corners, fill=blue!4, minimum height=10mm, minimum width=24mm, right=of in] (sys) {LTI system $h[n]$};
\node[right=14mm of sys] (out) {$H(e^{\jj\omega})\,e^{\jj\omega n}$};
\draw[->] (in) -- (sys);
\draw[->] (sys) -- (out);
\node[below=5mm of sys, gray!60!black, align=center, font=\small]
  {same frequency out,\\ scaled and phase-shifted by one number};
\end{tikzpicture}
\caption{The defining experiment. A pure complex exponential is the one input an
LTI system cannot reshape: it returns the same exponential, multiplied by the
single complex number $H(e^{\jj\omega})$.}
\label{ch14:fig-eigen}
\end{figure}

\begin{happeningbox}
Complex exponentials are the \emph{eigenfunctions} of every LTI system. Feed in
$e^{\jj\omega n}$ and the system cannot change its shape --- it is still a pure
rotation at frequency $\omega$. All the system is allowed to do is scale it and
spin its starting phase, and both of those are captured by the single complex
number $H(e^{\jj\omega})$. The function $H(e^{\jj\omega})$, swept over all
$\omega$, is the \emph{frequency response}, and it is just the eigenvalue
attached to each frequency.
\end{happeningbox}

\begin{connectionbox}
This is exactly the eigenvalue/eigenvector story from the linear algebra book,
lifted to infinite dimensions. Convolution by $h$ is a linear operator on the
space of sequences. The complex exponentials $e^{\jj\omega n}$ are its
eigenvectors, and $H(e^{\jj\omega})$ is the eigenvalue belonging to the
eigenvector of frequency $\omega$. Filtering is ``diagonal'' in the frequency
basis --- it just multiplies each frequency component by its own number. This
diagonalization provides the primary geometric justification for adopting the
Fourier basis: in these coordinates, every shift-invariant operation uncouples
into independent scalings.
\end{connectionbox}

\subsection{The frequency response is the DTFT of \texorpdfstring{$h$}{h}}

Look again at Equation~\eqref{ch14:Hdef}. It is the discrete-time Fourier
transform (DTFT) of the impulse response, from Chapter~8. So we can state the
central identity of this chapter in one line:

\begin{formalbox}
The frequency response is the DTFT of the impulse response:
\[
  H(e^{\jj\omega}) = \FT\{h[n]\} = \sum_{n=-\infty}^{\infty} h[n]\, e^{-\jj\omega n}.
\]
Conversely $h$ is the inverse DTFT of $H$,
\[
  h[n] = \frac{1}{2\pi}\int_{-\pi}^{\pi} H(e^{\jj\omega})\, e^{\jj\omega n}\,\dd\omega .
\]
Time and frequency descriptions of an LTI system carry exactly the same
information; the DTFT is the dictionary that translates one into the other.
\end{formalbox}

The notation $H(e^{\jj\omega})$ looks heavy, but it is deliberate. It reminds
you that $H$ is really a function of the point $e^{\jj\omega}$ on the unit
circle of the complex plane. As $\omega$ runs from $-\pi$ to $\pi$, the point
$e^{\jj\omega}$ walks once around that circle. Frequency response is what you
see as you stroll around the unit circle. (In Chapter~15 the $z$-transform will
let $z$ roam the whole plane; the frequency response is the slice of the
$z$-transform taken right on the circle $\abs{z}=1$.)

\begin{whybox}
Why does the convolution theorem fall out for free here? Because
Equation~\eqref{ch14:eigenrelation} says each eigen-input is scaled by
$H(e^{\jj\omega})$. An arbitrary input $x[n]$ is, by the inverse DTFT, a sum
(integral) of such eigen-inputs, one per frequency, with weights
$X(e^{\jj\omega})$. The system scales each one by its eigenvalue, so the output
weights are $Y(e^{\jj\omega}) = H(e^{\jj\omega})\,X(e^{\jj\omega})$.
Convolution in time becomes \emph{multiplication} in frequency. We did the work
in Chapter~11; the eigenfunction picture is why it had to be true.
\end{whybox}

\begin{historybox}
The idea that an entire dynamical system could be characterized by kicking it with an instantaneous,
infinitely narrow spike originated in telegraph engineering of the late 19th century.
Oliver Heaviside, an eccentric, largely self-taught English telegraph engineer working in the 1880s
and 1890s, worked on the practical problem of signal dispersion along transatlantic undersea
cables. Standard differential equations were slow to solve by hand for complex inputs.
Heaviside introduced his ``operational calculus,'' treating the derivative operator
$\dd/\dd t$ as an algebraic variable $p$, and introduced the unit step function. By differentiating
the step, he effectively created the impulse $\delta(t)$ decades before Paul Dirac formalized it in 1927.

Heaviside showed that the response to an impulse completely dictates the response to any input via
superposition. Pure mathematicians of the Royal Society rejected Heaviside's methods as
mathematically illegitimate: an operator is not an algebraic quantity, and a function that is zero everywhere
except at the origin where it is infinite with unit area violated classical analysis. Heaviside famously
responded: ``Shall I refuse my dinner because I do not fully understand the process of digestion?''
It took half a century for Laurent Schwartz (1950) to build the theory of distributions that vindicated
Heaviside and Dirac, while George Campbell and Otto Zobel at AT\&T (1910s--1920s) translated operational
ideas into the modern frequency-selective filters that made nationwide telephony possible.

This unified the two views: any stable linear system is completely characterized either by its
impulse response $h[n]$ in time, or by its frequency response $H(e^{\jj\omega})$ across the spectrum.
Yet the discrete-time frequency response has a strict limitation: it lives solely on the unit circle $|z|=1$,
and it converges only for systems that are strictly stable. If a system has poles outside the circle, or if
we want to study unbounded growth and feedback stability, the Fourier sum diverges. Chapter~15
removes this limitation by introducing the $z$-transform, using the complex plane to analyze stability,
causality, and frequency response in a single geometric framework.
\end{historybox}

\begin{focusbox}
Two facts to carry through the rest of the book:
\begin{enumerate}
  \item $H(e^{\jj\omega})$ is periodic in $\omega$ with period $2\pi$. A
  discrete-time system only ``knows'' frequencies on one trip around the unit
  circle; everything else is an alias.
  \item If $h[n]$ is real (as it is for every filter we will actually build),
  then $H(e^{-\jj\omega}) = \overline{H(e^{\jj\omega})}$. The magnitude is even
  and the phase is odd, so we usually plot only $0 \le \omega \le \pi$.
\end{enumerate}
\end{focusbox}

\subsection{Impulse response versus step response}
\label{ch14:step-response}

While the impulse response $h[n]$ theoretically characterizes any LTI system, injecting a true unit impulse into physical hardware is practically challenging: an impulse requires high peak amplitude over a single sample that can saturate sensor amplifiers or ADCs. In practice, engineers frequently test systems by applying a \textbf{unit step input}:
\[
  u[n] = \begin{cases} 1, & n \ge 0, \\ 0, & n < 0. \end{cases}
\]
Because the step input is the running accumulation of unit impulses, $u[n] = \sum_{k=0}^\infty \delta[n-k]$, linearity and time-invariance guarantee that the \textbf{step response} $s[n] = T\{u[n]\}$ is the running sum of the impulse response:
\begin{equation}
  s[n] = \sum_{k=-\infty}^{n} h[k], \qquad h[n] = s[n] - s[n-1].
  \label{ch14:step-impulse-rel}
\end{equation}
The two responses describe the same dynamics from two angles:
\begin{itemize}
  \item \textbf{Impulse response $h[n]$}: shows instantaneous memory decay, individual echoes, and resonant oscillation frequencies.
  \item \textbf{Step response $s[n]$}: reveals the steady-state DC gain $\lim_{n\to\infty} s[n] = H(e^{\jj 0}) = \sum_k h[k]$, rise time (how rapidly the system responds to a sudden level shift), and transient overshoot before settling.
\end{itemize}

\begin{figure}[H]
\centering
\begin{tikzpicture}
\begin{axis}[
title style={font=\small},
  name=p1,
  width=6.6cm, height=4.6cm,
  title={Impulse response $h[n]$ (transient)},
  title style={font=\small},
  xlabel={$n$}, ylabel={$h[n]$},
  xmin=-1, xmax=20, ymin=-0.45, ymax=1.15,
  xtick={0,5,10,15,20}, ytick={-0.2,0,0.5,1.0},
  axis lines=left, ymajorgrids, grid style={gray!20},
]
\addplot+[ycomb, mark=*, sigblue, thick, mark options={fill=sigblue}]
  coordinates {
    (0,1.0)(1,0.54)(2,0.05)(3,-0.25)(4,-0.32)(5,-0.24)(6,-0.10)(7,0.01)
    (8,0.08)(9,0.09)(10,0.06)(11,0.02)(12,-0.01)(13,-0.02)(14,-0.02)
    (15,-0.01)(16,0.0)(17,0.0)(18,0.0)(19,0.0)(20,0.0)
  };
\node[above, font=\scriptsize, sigblue] at (axis cs:0,1.02) {initial tap};
\end{axis}

\begin{axis}[
title style={font=\small},
  name=p2, at={(p1.right of south east)}, anchor=left of south west, xshift=1.2cm,
  width=6.6cm, height=4.6cm,
  title={Step response $s[n] = \sum_{k=0}^n h[k]$},
  title style={font=\small},
  xlabel={$n$}, ylabel={$s[n]$},
  xmin=-1, xmax=20, ymin=0, ymax=1.4,
  xtick={0,5,10,15,20}, ytick={0,0.5,0.9,1.0,1.2},
  axis lines=left, ymajorgrids, grid style={gray!20},
]
\addplot+[ycomb, mark=square*, specred, thick, mark options={fill=specred}]
  coordinates {
    (0,1.0)(1,1.24)(2,1.29)(3,1.04)(4,0.72)(5,0.48)(6,0.38)(7,0.39)
    (8,0.47)(9,0.56)(10,0.62)(11,0.64)(12,0.63)(13,0.61)(14,0.59)
    (15,0.58)(16,0.58)(17,0.58)(18,0.58)(19,0.58)(20,0.58)
  };
\addplot[dashed, gray!70, thick] coordinates {(-1,0.58)(20,0.58)};
\node[above, font=\scriptsize, gray!80!black] at (axis cs:10,0.62) {DC gain $H(1) = 0.58$};
\node[above, font=\scriptsize, specred] at (axis cs:2,1.31) {overshoot};
\end{axis}
\end{tikzpicture}
\caption{Impulse response versus step response for an underdamped resonant system. Left: the impulse response $h[n]$ rings and decays toward zero. Right: the step response $s[n] = \sum_{k=0}^n h[k]$ integrates the ring, exhibiting overshoot before settling at the DC gain $H(e^{\jj 0}) = \sum_k h[k] = 0.58$.}
\label{ch14:fig-impulse-vs-step}
\end{figure}

\section{Magnitude and phase: what the complex number means}
\label{ch14:magphase}

$H(e^{\jj\omega})$ is complex, so like any complex number it has a magnitude and
an angle:
\[
  H(e^{\jj\omega}) = \abs{H(e^{\jj\omega})}\; e^{\jj\,\angle H(e^{\jj\omega})}.
\]
Put this back into the eigen-relation~\eqref{ch14:eigenrelation} and write the
input as a real cosine, $x[n] = \cos(\omega n)$. By linearity and the fact that
$\cos(\omega n) = \tfrac12(e^{\jj\omega n} + e^{-\jj\omega n})$, the steady-state
output of a real system is
\begin{equation}
  x[n] = \cos(\omega n) \;\;\longrightarrow\;\;
  y[n] = \abs{H(e^{\jj\omega})}\,\cos\!\big(\omega n + \angle H(e^{\jj\omega})\big).
  \label{ch14:cosinout}
\end{equation}

So the two parts of the complex number have crisp physical meanings.

\begin{itemize}
  \item \textbf{Magnitude response} $\abs{H(e^{\jj\omega})}$: the \emph{gain} at
  frequency $\omega$. A value of $2$ doubles a sine of that frequency; $0.1$
  knocks it down to a tenth; $0$ kills it outright. This is the curve people
  usually mean when they say ``the filter's response.''
  \item \textbf{Phase response} $\angle H(e^{\jj\omega})$: the \emph{phase shift}
  the filter adds at frequency $\omega$. A shift of $-\pi/2$ turns a cosine into
  a sine; in time it slides that sine later (a delay).
\end{itemize}

\begin{intuitionbox}
A cosine goes in, a cosine of the \emph{same frequency} comes out --- never a
different frequency. An LTI system cannot create new tones; it can only turn the
volume knob (magnitude) and slide the start time (phase) of each tone already
present. That is all a linear filter can do. Distortion that
generates harmonics --- the growl of an overdriven guitar amp --- is exactly the
signature of \emph{non}-linearity.
\end{intuitionbox}

\subsection{Decibels: why gain is plotted on a log scale}

Magnitude responses span enormous ranges --- a good filter might pass a band at
gain $1$ and reject another at gain $0.001$. To see both at once we use
decibels:
\[
  \abs{H}_{\text{dB}} = 20\log_{10}\abs{H(e^{\jj\omega})}.
\]
A gain of $1$ is $0$~dB; gain $\tfrac{1}{\sqrt2}\approx 0.707$ is $-3$~dB; gain
$0.1$ is $-20$~dB; gain $0.001$ is $-60$~dB. The factor $20$ (not $10$) is
because power goes as amplitude squared and
$10\log_{10}(\abs{H}^2)=20\log_{10}\abs{H}$. The famous ``$-3$~dB point'' marks
where a filter has dropped to half \emph{power}, and it is the conventional edge
of a passband.

\begin{trapbox}
Do not confuse the unit-circle frequency $\omega$ (radians per sample, range
$-\pi$ to $\pi$) with an ordinary frequency in hertz. If the sampling rate is
$f_s$, then $\omega = 2\pi f / f_s$. The top of the digital band, $\omega=\pi$,
is the Nyquist frequency $f_s/2$. We return to this in Chapter~18; for now,
just remember that ``$\omega = \pi$'' means ``the fastest wiggle a sampled
signal can represent: flipping sign every sample.''
\end{trapbox}

\section{A first worked example: the two-tap averager}
\label{ch14:avg}

The smallest real filter makes this concrete. Take the
two-point moving average,
\[
  y[n] = \tfrac12\,x[n] + \tfrac12\,x[n-1].
\]
Its impulse response is read straight off the coefficients: poke it with
$\delta[n]$ and you get $h[0]=\tfrac12$, $h[1]=\tfrac12$, all else zero. Plug
into Equation~\eqref{ch14:Hdef}:
\[
  H(e^{\jj\omega}) = \tfrac12 + \tfrac12 e^{-\jj\omega}.
\]
Now do the standard trick: factor out half the total phase to expose a clean
cosine. Pull out $e^{-\jj\omega/2}$:
\[
  H(e^{\jj\omega})
   = e^{-\jj\omega/2}\left(\tfrac12 e^{\jj\omega/2} + \tfrac12 e^{-\jj\omega/2}\right)
   = e^{-\jj\omega/2}\,\cos\!\big(\tfrac{\omega}{2}\big).
\]
This form is convenient because it gives magnitude and phase separately:
\[
  \abs{H(e^{\jj\omega})} = \abs{\cos\!\big(\tfrac{\omega}{2}\big)},
  \qquad
  \angle H(e^{\jj\omega}) = -\tfrac{\omega}{2}.
\]

At DC ($\omega=0$) the gain is $\cos 0 = 1$: a constant passes
untouched, which makes sense, the average of a constant is itself. At Nyquist
($\omega=\pi$) the gain is $\cos(\pi/2)=0$: the alternating sequence
$\dots,1,-1,1,-1,\dots$ averages adjacent samples to zero and is annihilated.
In between, the gain falls smoothly. This is a (gentle) \emph{low-pass} filter:
it favors slow wiggles and suppresses fast ones. And the phase is a perfectly
straight line, $-\omega/2$. We will see in a moment that a straight-line phase
is the best kind to have.

\begin{figure}[H]
\centering
\begin{tikzpicture}
\begin{axis}[
  width=11.5cm, height=6.2cm,
  xlabel={$\omega$ (radians/sample)}, ylabel={magnitude $\abs{H}$},
  domain=0:pi, samples=120,
  xmin=0, xmax=pi, ymin=0, ymax=1.05,
  xtick={0,1.5708,3.1416}, xticklabels={$0$,$\pi/2$,$\pi$},
  ytick={0,0.5,0.707,1}, yticklabels={$0$,$0.5$,$0.707$,$1$},
  grid=both, grid style={gray!20},
  legend pos=south west, legend cell align=left,
  axis lines=box,
]
\addplot[sigblue, very thick] {abs(cos(deg(x/2)))};
\addlegendentry{two-tap averager (low-pass)}
\addplot[specred, very thick, dashed] {abs(sin(deg(x/2)))};
\addlegendentry{two-tap difference (high-pass)}
\addplot[gray, dotted, thick] coordinates {(0,0.707)(3.1416,0.707)};
\end{axis}
\end{tikzpicture}
\caption{Magnitude responses of the two simplest filters. The averager
$h=[\tfrac12,\tfrac12]$ passes low frequencies and nulls Nyquist; the
difference $h=[\tfrac12,-\tfrac12]$ does the opposite. The dotted line marks the
$-3$~dB ($0.707$) half-power level.}
\label{ch14:fig-avg}
\end{figure}

\subsection{The mirror image: a two-tap difference}

Flip one sign: $y[n] = \tfrac12 x[n] - \tfrac12 x[n-1]$, so $h=[\tfrac12,-\tfrac12]$.
The same algebra gives
\[
  H(e^{\jj\omega}) = \tfrac12 - \tfrac12 e^{-\jj\omega}
   = \jj\,e^{-\jj\omega/2}\,\sin\!\big(\tfrac{\omega}{2}\big),
  \qquad
  \abs{H(e^{\jj\omega})} = \abs{\sin\!\big(\tfrac{\omega}{2}\big)}.
\]
Now the gain is $0$ at DC and rises to $1$ at Nyquist --- a \emph{high-pass}
filter (the dashed curve in Figure~\ref{ch14:fig-avg}). The intuition is just as
clean: a difference asks ``how much did the signal change between samples?'' A
constant does not change, so it is removed; a fast alternation changes maximally,
so it survives. First-differencing is the discrete-time counterpart of differentiation,
and like differentiation it emphasizes high frequencies.

\begin{example}[Predicting the output by hand]
\label{ch14:ex-handpredict}
Send $x[n] = \cos(\tfrac{\pi}{2} n)$ through the averager. Here $\omega=\pi/2$,
so $\abs{H}=\cos(\pi/4)=0.707$ and $\angle H = -\pi/4$. By
Equation~\eqref{ch14:cosinout} the steady-state output is
\[
  y[n] = 0.707\,\cos\!\big(\tfrac{\pi}{2}n - \tfrac{\pi}{4}\big).
\]
Check it directly: $x = \dots,1,0,-1,0,1,\dots$ and
$y[n]=\tfrac12(x[n]+x[n-1]) = \dots, 0.5, 0.5, -0.5, -0.5, \dots$, which is
indeed a cosine of amplitude $0.707$ shifted by an eighth of a period. The
frequency response told us the answer without grinding through the convolution.
\end{example}

\section{Group delay: the slope of the phase}
\label{ch14:groupdelay}

Magnitude is intuitive; phase takes more care. The raw phase
$\angle H(e^{\jj\omega})$ is hard to read because a delay shows up as a phase
that grows with frequency, and because phase wraps around every $2\pi$. The
quantity that turns phase into something physical is the \emph{group delay}:
\begin{equation}
  \tau_g(\omega) = -\frac{\dd}{\dd\omega}\,\angle H(e^{\jj\omega}).
  \label{ch14:gd}
\end{equation}
Group delay measures, in samples, how long the filter holds back the
\emph{envelope} of a narrowband signal centered at $\omega$. If you send in a
short pulse whose energy sits near frequency $\omega$, the bump comes out
$\tau_g(\omega)$ samples later.

\begin{whybox}
Why the \emph{derivative} of phase, and why does it act on the envelope? Picture
a narrowband packet built from frequencies $\omega_0 + \delta\omega$ clustered tightly around a carrier $\omega_0$. Near $\omega_0$, approximate the phase curve by its tangent line, $\angle H(e^{\jj(\omega_0+\delta\omega)}) \approx \phi_0 - \tau_g\,\delta\omega$. Each component exponential $e^{\jj(\omega_0+\delta\omega)n}$ then emerges with phase
\[
  e^{\jj(\omega_0+\delta\omega)n}\,e^{\jj(\phi_0 - \tau_g\,\delta\omega)}
  \;=\; \underbrace{e^{\jj(\omega_0 n + \phi_0)}}_{\text{shifted carrier}}\;
        \underbrace{e^{\jj\delta\omega\,(n - \tau_g)}}_{\text{envelope delayed by }\tau_g}.
\]
The fast carrier shifts by $-\phi_0/\omega_0$ (\emph{phase delay}), whereas all the slow modulation offsets $\delta\omega$ line up at time $n - \tau_g$, so the packet's envelope reappears intact $\tau_g = -\dd(\angle H)/\dd\omega$ samples later.
\end{whybox}

\begin{intuitionbox}
\textbf{Linear phase = constant group delay = no waveform distortion.} If
$\angle H(e^{\jj\omega}) = -\alpha\omega$ for some constant $\alpha$, then
$\tau_g = \alpha$ for \emph{every} frequency. All frequencies are delayed by the
same $\alpha$ samples, so the output is just the input shifted in time (and
rescaled in amplitude band by band). The shape of the signal is preserved.

If instead the slope varies with $\omega$, different frequencies arrive at
different times. Bass lags treble, the waveform smears, and a crisp impulse
turns into a ringing blur. That smearing is \emph{phase distortion}, and it is
inaudible in some applications and ruinous in others.
\end{intuitionbox}

Our two-tap averager had $\angle H = -\omega/2$, a perfectly straight line, so
$\tau_g = \tfrac12$ sample at every frequency. It is a linear-phase filter ---
the whole signal is delayed by half a sample and nothing is smeared. This is no
accident: any FIR filter whose taps are symmetric (or antisymmetric) about their
center has exactly linear phase, a property we will exploit heavily in
Chapter~16.

\begin{figure}[H]
\centering
\begin{tikzpicture}
\begin{axis}[
  width=11.5cm, height=5.2cm,
  xlabel={$\omega$}, ylabel={phase $\angle H$ (rad)},
  domain=0:pi, samples=2,
  xmin=0, xmax=pi, ymin=-1.7, ymax=0.1,
  xtick={0,1.5708,3.1416}, xticklabels={$0$,$\pi/2$,$\pi$},
  grid=both, grid style={gray!20}, axis lines=box,
  legend pos=north east, legend cell align=left,
  legend style={fill=white, font=\footnotesize},
]
\addplot[sigblue, very thick] {-x/2};
\addlegendentry{linear phase ($\tau_g=0.5$, no smear)}
\addplot[specred, very thick, dashed, domain=0:pi, samples=80] {-0.5*x - 0.45*sin(deg(2*x))};
\addlegendentry{nonlinear phase (frequency-dependent delay)}
\end{axis}
\end{tikzpicture}
\caption{A straight phase line delays all frequencies equally and preserves
waveform shape. A curved phase line delays them differently, smearing transients
even when the magnitude response is identical.}
\label{ch14:fig-phase}
\end{figure}

\begin{trapbox}
Two filters can have \emph{identical} magnitude responses and behave completely
differently, because their phases differ. ``It has a flat passband'' tells you
nothing about whether it will ring or smear. Always check the phase (or group
delay) when transient shape matters --- and it almost always matters for the
short, sharp motions of a gesture.
\end{trapbox}

\section{Reading a frequency response like a map}
\label{ch14:reading}

Hand an engineer a magnitude plot and they read it the way you read a face. Here
is the vocabulary.

\begin{itemize}
  \item \textbf{Passband:} the range of $\omega$ where the gain is near its
  maximum (ideally near $0$~dB). These frequencies survive.
  \item \textbf{Stopband:} the range where the gain is small (deeply negative
  dB). These frequencies are rejected.
  \item \textbf{Transition band:} the slope between them. A steep transition is
  ``sharp''; a gentle one is ``soft.'' Sharpness costs taps (longer $h$) and
  costs delay.
  \item \textbf{Cutoff frequency} $\omega_c$: by convention the $-3$~dB point,
  where the gain has fallen to $0.707$ (half power).
  \item \textbf{Ripple:} small wiggles in the passband or stopband. Some filter
  designs trade passband flatness for a steeper transition; the wiggles are the
  price.
  \item \textbf{Notch / null:} a frequency where the gain is exactly (or nearly)
  zero, which removes a single tone such as 60~Hz mains hum.
  \item \textbf{Resonant peak:} a frequency where the gain spikes \emph{above}
  the passband --- the system rings there. More on this in
  Section~\ref{ch14:resonance}.
\end{itemize}

\begin{figure}[H]
\centering
\begin{tikzpicture}
\begin{axis}[
  width=12.5cm, height=6cm,
  xlabel={$\omega$}, ylabel={$\abs{H}$},
  domain=0:pi, samples=160,
  xmin=0, xmax=pi, ymin=0, ymax=1.15,
  xtick={0,1.0472,3.1416}, xticklabels={$0$,$\omega_c$,$\pi$},
  ytick={0,0.707,1}, yticklabels={$0$,$0.707$,$1$},
  grid=both, grid style={gray!20}, axis lines=box, clip=false,
]
% a smooth low-pass-ish curve
\addplot[sigblue, very thick] {1/sqrt(1 + (x/1.0472)^8)};
\node[anchor=north west, sigblue] at (axis cs:0.15,0.85) {\small passband};
\node[anchor=south, specred] at (axis cs:2.55,0.06) {\small stopband};
\node[anchor=west] at (axis cs:1.07,0.5) {\small transition};
\draw[gray,dashed] (axis cs:1.0472,0) -- (axis cs:1.0472,0.707);
\draw[gray,dotted] (axis cs:0,0.707) -- (axis cs:1.0472,0.707);
\node[anchor=west, gray!50!black] at (axis cs:1.15,0.72) {\small $-3$~dB};
\end{axis}
\end{tikzpicture}
\caption{Anatomy of a low-pass magnitude response. The cutoff $\omega_c$ is read
at the $-3$~dB (half-power, $0.707$) crossing.}
\label{ch14:fig-anatomy}
\end{figure}

\begin{walkbox}
Here is a magnitude curve that is $\approx 1$ from $\omega=0$ up to $\omega=0.4\pi$,
then drops fast to near $0$ and stays there. What does this filter do to a
recording of speech sampled at 8~kHz?

The Nyquist frequency is $4$~kHz, mapped to $\omega=\pi$. So $\omega=0.4\pi$ is
$0.4\times 4 = 1.6$~kHz. The filter passes everything below about 1.6~kHz and
removes everything above. Speech keeps its vowels and low formants but loses the
crisp ``s'' and ``f'' sounds, which live higher up. It will sound muffled, as if
heard through a wall. You just read a filter without computing a single
convolution.
\end{walkbox}

\section{How a system reshapes a spectrum}
\label{ch14:reshape}

We now have everything we need to picture filtering end to end. By the
convolution theorem, in the frequency domain the system simply multiplies the
input spectrum by its frequency response, point by point:
\begin{equation}
  Y(e^{\jj\omega}) = H(e^{\jj\omega})\,X(e^{\jj\omega}).
  \label{ch14:multiply}
\end{equation}
And magnitudes multiply while phases add:
\[
  \abs{Y} = \abs{H}\cdot\abs{X}, \qquad \angle Y = \angle H + \angle X.
\]

\begin{happeningbox}
Filtering is \emph{spectral shaping}. Lay the input spectrum $\abs{X}$ over the
filter's gain curve $\abs{H}$. Where the curve is high, those frequencies of the
input pass through loud; where it dips, they are turned down; where it hits zero,
they vanish. The output spectrum is the input spectrum ``painted'' by the
filter's gain. Nothing new appears; the existing colors are just re-weighted.
\end{happeningbox}

\begin{figure}[H]
\centering
\begin{tikzpicture}
\begin{axis}[
  width=12.5cm, height=5.6cm,
  xlabel={$\omega$}, ylabel={magnitude},
  domain=0:pi, samples=160,
  xmin=0, xmax=pi, ymin=0, ymax=1.3,
  xtick={0,1.5708,3.1416}, xticklabels={$0$,$\pi/2$,$\pi$},
  grid=both, grid style={gray!20}, axis lines=box,
  legend pos=north east, legend cell align=left,
]
% input: two bumps, a low one and a high one
\addplot[gray, thick] {0.9*exp(-((x-0.6)*3.0)^2) + 0.8*exp(-((x-2.4)*3.0)^2)};
\addlegendentry{input $\abs{X}$ (two tones)}
% filter: low-pass gain
\addplot[sigblue, very thick, dashed] {1/sqrt(1+(x/1.0)^6)};
\addlegendentry{filter $\abs{H}$ (low-pass)}
% output = product
\addplot[specred, very thick] {(0.9*exp(-((x-0.6)*3.0)^2) + 0.8*exp(-((x-2.4)*3.0)^2)) * 1/sqrt(1+(x/1.0)^6)};
\addlegendentry{output $\abs{Y}=\abs{H}\abs{X}$}
\end{axis}
\end{tikzpicture}
\caption{An input with a low-frequency tone and a high-frequency tone, sent
through a low-pass filter. The low bump survives nearly intact; the high bump is
crushed where $\abs{H}$ is small. Output magnitude is the pointwise product.}
\label{ch14:fig-reshape}
\end{figure}

\begin{example}[Killing 60~Hz hum]
\label{ch14:ex-notch}
A sensor sampled at $f_s = 1000$~Hz picks up mains interference at $60$~Hz. In
digital terms that tone sits at $\omega_0 = 2\pi\cdot 60/1000 = 0.377$~rad. To
remove it we want a frequency response that is $\approx 1$ everywhere except a
deep null at $\pm\omega_0$. A two-zero ``notch'' does it: place a transmission
zero exactly at $e^{\pm\jj\omega_0}$ so that $H(e^{\jj\omega_0})=0$. We will
build the $h[n]$ for this from its zeros in Chapter~15, but the point now is
that we specify the filter by drawing the response we want --- a flat line with
one narrow notch --- and then work backward to the coefficients.
\end{example}

\section{Resonance: when a system rings}
\label{ch14:resonance}

So far our filters only attenuated. But many systems amplify a narrow band of
frequencies dramatically --- a wine glass, a guitar body, a tuned circuit, a
mass on a spring. In the frequency response this shows up as a tall, narrow peak.
In the time domain it shows up as \emph{ringing}: poke the system once and it
keeps oscillating for a while, decaying slowly. These are two views of the same
thing.

The simplest resonator is a single-pole feedback filter,
\begin{equation}
  y[n] = x[n] + r\,e^{\jj\omega_0}\, y[n-1] \quad\text{(complex)},
  \qquad 0 < r < 1,
  \label{ch14:reson}
\end{equation}
or, in real form using a conjugate pair of poles,
\[
  y[n] = x[n] + 2r\cos(\omega_0)\,y[n-1] - r^2\,y[n-2].
\]
Its impulse response is a decaying sinusoid, $h[n] \approx r^n\cos(\omega_0 n)$
for $n\ge0$: an exponentially fading ring at frequency $\omega_0$. The closer
$r$ creeps toward $1$, the slower the decay and --- this is the duality --- the
\emph{taller and narrower} the spectral peak.

\begin{intuitionbox}
A resonance is feedback that nearly --- but not quite --- sustains itself at one
frequency. Energy fed in near $\omega_0$ gets recycled around the feedback loop
many times before it leaks away, building up a large response. The number of
times it recycles is roughly the \emph{quality factor} $Q$: high $Q$ means a
tall sharp peak in frequency and a long ring in time; low $Q$ means a broad
bump that dies quickly. ``Sharp in frequency'' and ``long in time'' are the same
statement --- another instance of the time--frequency uncertainty tradeoff.
\end{intuitionbox}

\begin{figure}[H]
\centering
\begin{tikzpicture}
\begin{axis}[
title style={font=\small},
  width=6.1cm, height=5cm,
  title={Impulse response $h[n]$},
  title style={font=\small},
  xlabel={$n$}, ylabel={$h[n]$},
  xmin=0, xmax=30, ymin=-1.05, ymax=1.05,
  ytick={-1,0,1}, grid=both, grid style={gray!20}, axis lines=box,
  samples=31, domain=0:30,
]
\addplot[ycomb, sigblue, thick, mark=*, mark size=1pt]
  {0.9^x * cos(deg(0.9*x))};
\end{axis}
\begin{axis}[
  at={(7.4cm,0)},
  width=6.1cm, height=5cm,
  title={Magnitude $|H(e^{\jj\omega})|$},
  title style={font=\small},
  xlabel={$\omega$}, ylabel={$\abs{H}$},
  domain=0:pi, samples=160,
  xmin=0, xmax=pi, ymin=0, ymax=6,
  xtick={0,0.9,3.1416}, xticklabels={$0$,$\omega_0$,$\pi$},
  grid=both, grid style={gray!20}, axis lines=box,
]
\addplot[specred, very thick]
  {1/sqrt((1 - 2*0.9*cos(deg(x-0.9)) + 0.9^2)*(1 - 2*0.9*cos(deg(x+0.9)) + 0.9^2))};
\end{axis}
\end{tikzpicture}
\caption{A resonator with poles at radius $r=0.9$ and angle $\omega_0=0.9$~rad.
Left: the impulse response is a decaying ring. Right: the matching magnitude
response has a tall narrow peak at $\omega_0$. Long ring $\Leftrightarrow$
sharp peak.}
\label{ch14:fig-reson}
\end{figure}

\begin{whybox}
Where does the peak come from, geometrically? Write the frequency response as a
ratio of distances in the complex plane (Chapter~15 makes this precise). The
denominator contains a factor $\abs{e^{\jj\omega} - p}$, the distance from the
unit-circle point $e^{\jj\omega}$ to the pole $p = r\,e^{\jj\omega_0}$. As
$\omega$ sweeps past $\omega_0$, the unit-circle point glides closest to the
pole, that distance shrinks, and the response --- being one over the distance ---
shoots up. The nearer the pole sits to the circle ($r\to 1$), the smaller the
minimum distance, and the higher the spike. Resonance is just ``the unit circle
passing close to a pole'' (Figure~\ref{ch14:fig-polezero}).
\end{whybox}

\begin{figure}[H]
\centering
\begin{tikzpicture}[scale=2.4]
% unit circle
\draw[gray!55, thick] (0,0) circle (1);
% axes
\draw[->, gray!70] (-1.25,0) -- (1.3,0) node[right, black] {\small Re};
\draw[->, gray!70] (0,-1.25) -- (0,1.3) node[above, black] {\small Im};
% pole angle line
\draw[gray!50, dashed] (0,0) -- ({0.9*cos(51.57)},{0.9*sin(51.57)});
% conjugate poles at r=0.9, angle 0.9 rad ~ 51.57 deg
\node[specred] at ({0.9*cos(51.57)},{0.9*sin(51.57)}) {$\times$};
\node[specred] at ({0.9*cos(-51.57)},{0.9*sin(-51.57)}) {$\times$};
\node[specred, anchor=north west] at ({0.9*cos(51.57)+0.02},{0.9*sin(51.57)-0.02}) {\small $p=re^{\jj\omega_0}$};
% the unit-circle test point near the pole
\fill[sigblue] ({cos(51.57)},{sin(51.57)}) circle (0.018);
\node[sigblue, anchor=south west] at ({cos(51.57)+0.02},{sin(51.57)+0.02}) {\small $e^{\jj\omega}$};
% short-distance segment
\draw[sigblue, thick] ({cos(51.57)},{sin(51.57)}) -- ({0.9*cos(51.57)},{0.9*sin(51.57)});
\node[sigblue, anchor=south, rotate=51.57] at ({0.95*cos(51.57)},{0.95*sin(51.57)}) {};
\end{tikzpicture}
\caption{Pole-zero view of the resonator. The conjugate poles ($\times$) sit just
inside the unit circle at radius $r=0.9$ and angle $\pm\omega_0$. As the test
point $e^{\jj\omega}$ slides around the circle to $\omega_0$, its distance to the
pole (blue segment) reaches a minimum, so the response --- one over that distance
--- peaks. Push the pole toward the circle and the peak grows taller and
narrower.}
\label{ch14:fig-polezero}
\end{figure}

\begin{applicationbox}
Resonance is everywhere in real instruments. The body of an acoustic guitar has
strong resonances that color its tone; a vocal tract has resonances (formants)
that distinguish ``ah'' from ``ee''; a snare drum has a tuned ring. When we
model or synthesize these sounds (Chapters~28--29), we build banks of resonators
and set each pole's angle (pitch) and radius (decay time). Conversely, to
\emph{analyze} a sound we look for spectral peaks and infer the resonances that
produced them.
\end{applicationbox}

\section{Putting numbers in: a worked FIR low-pass}
\label{ch14:firworked}

Let us compute a frequency response we cannot do in our heads, to see the
machinery turn. Take a length-5 averaging filter,
\[
  h = \tfrac15\,[\,1,\,1,\,1,\,1,\,1\,], \qquad h[0..4]=0.2.
\]
By Equation~\eqref{ch14:Hdef},
\[
  H(e^{\jj\omega}) = \tfrac15\sum_{n=0}^{4} e^{-\jj\omega n}.
\]
This is a finite geometric series with ratio $e^{-\jj\omega}$. Summing the five terms and factoring out the half-angle phases from top and bottom (the Dirichlet-kernel trick) gives
\[
  H(e^{\jj\omega})
  = \frac{1}{5}\,\frac{1 - e^{-\jj 5\omega}}{1 - e^{-\jj\omega}}
  = \frac{1}{5}\,\frac{e^{-\jj 5\omega/2}\bigl(e^{\jj 5\omega/2} - e^{-\jj 5\omega/2}\bigr)}
                      {e^{-\jj\omega/2}\bigl(e^{\jj\omega/2} - e^{-\jj\omega/2}\bigr)}
  = e^{-\jj 2\omega}\,\frac{1}{5}\,\frac{\sin(5\omega/2)}{\sin(\omega/2)}.
\]
The magnitude is the (aliased) sinc shape
$\abs{H} = \tfrac15\abs{\sin(5\omega/2)/\sin(\omega/2)}$, and the phase is the
straight line $-2\omega$ (so $\tau_g = 2$ samples --- the center of a 5-tap
symmetric filter, exactly as the linear-phase rule predicts). It is low-pass,
with a main lobe near DC and decaying side lobes --- the longer the average, the
narrower the main lobe and the lower the cutoff.

\begin{figure}[H]
\centering
\begin{tikzpicture}
\begin{axis}[
  width=12cm, height=5.4cm,
  xlabel={$\omega$}, ylabel={$\abs{H}$ (length-5 box)},
  domain=0.001:pi, samples=180,
  xmin=0, xmax=pi, ymin=0, ymax=1.05,
  xtick={0,1.2566,2.5133,3.1416},
  xticklabels={$0$,$2\pi/5$,$4\pi/5$,$\pi$},
  grid=both, grid style={gray!20}, axis lines=box,
]
\addplot[sigblue, very thick] {abs( (sin(deg(2.5*x))) / (5*sin(deg(0.5*x))) )};
\node[anchor=south, sigblue] at (axis cs:0.25,0.9) {\small main lobe};
\node[anchor=south] at (axis cs:1.9,0.12) {\small side lobes};
\end{axis}
\end{tikzpicture}
\caption{Magnitude response of the length-5 box-average filter. Note the nulls
at $\omega = 2\pi/5$ and $4\pi/5$ (where $\sin(5\omega/2)=0$) and the leaky side
lobes --- the price of an abrupt, rectangular impulse response.}
\label{ch14:fig-box5}
\end{figure}

\begin{trapbox}
Those side lobes are why a plain moving average is a \emph{poor} low-pass filter:
its stopband attenuation is only about $-13$~dB at the first side lobe, so it
leaks high-frequency content. The cure is to taper the taps smoothly to zero
instead of chopping them off square --- this is windowing, the subject of
Chapter~16. The square impulse response causes the leaky spectrum; a smoother
$h$ buys a cleaner $H$.
\end{trapbox}

\begin{labbox}
Compute and plot frequency responses yourself. The most direct way is to evaluate
the DTFT sum in Equation~\eqref{ch14:Hdef} on a dense frequency grid; in practice we sample it densely with
an FFT (zero-padded), which is what \texttt{freqz} does.

\begin{lstlisting}[language=Python]
import numpy as np
import matplotlib.pyplot as plt

def freq_response(h, N=1024):
    w = np.linspace(0, np.pi, N)
    # H(e^{jw}) = sum_n h[n] e^{-j w n}, evaluated at each w
    n = np.arange(len(h))
    H = h @ np.exp(-1j * np.outer(n, w))   # shape (N,)
    return w, H

h = np.ones(5) / 5.0                        # length-5 box average
w, H = freq_response(h)

plt.subplot(2,1,1)
plt.plot(w, 20*np.log10(np.abs(H) + 1e-9))  # magnitude in dB
plt.ylabel("|H| (dB)"); plt.ylim(-60, 5)

plt.subplot(2,1,2)
plt.plot(w, np.unwrap(np.angle(H)))         # unwrap before reading slope
plt.xlabel("omega (rad/sample)"); plt.ylabel("phase (rad)")
plt.show()

# group delay = -d(phase)/d(omega), as a finite difference
phase = np.unwrap(np.angle(H))
gd = -np.gradient(phase, w)
print("group delay near DC:", gd[1])        # ~ 2 samples
\end{lstlisting}
Swap in \texttt{h = np.array([0.5, 0.5])} or \texttt{[0.5, -0.5]} and confirm the
curves in Figure~\ref{ch14:fig-avg}. Then try a longer box and watch the cutoff
move down. Always \texttt{unwrap} the phase before differentiating, or the
$2\pi$ jumps will wreck your group delay.
\end{labbox}

\section{The two views, side by side}
\label{ch14:summary-views}

It is worth collecting the dictionary in one place, because moving fluently
between the two views is the core skill.

\begin{center}
\renewcommand{\arraystretch}{1.3}
\begin{tabular}{@{}p{0.46\textwidth} p{0.46\textwidth}@{}}
\toprule
\textbf{Time domain: impulse response $h[n]$} & \textbf{Frequency domain: response $H(e^{\jj\omega})$} \\
\midrule
What the system does to one poke & What it does to each pure tone \\
Output via convolution $y=h\conv x$ & Output via multiplication $Y=HX$ \\
Length of $h$ = memory / latency & Sharpness of transitions \\
Symmetric taps & Linear phase, constant group delay \\
Long, slowly decaying ring & Tall, narrow resonant peak \\
Abrupt (boxy) $h$ & Leaky side lobes in $\abs{H}$ \\
Delay by $d$ samples: $h[n]=\delta[n-d]$ & Pure phase ramp $e^{-\jj\omega d}$, flat gain \\
\bottomrule
\end{tabular}
\end{center}

\begin{aestheticbox}
The main idea of this chapter is that \emph{a filter is one object with two
faces}. Designers think in frequency (``flat to 1~kHz, then 40~dB down'') and
implement in time (a list of taps run through a convolution loop on a chip).
Every property you care about --- latency, ringing, sharpness, distortion ---
can be read from either face once you know how the DTFT maps between them. Much of
the rest of the book builds skill at reading both faces at once.
\end{aestheticbox}

\begin{goalbox}
For the wrist gesture controller, the frequency response is the lens through
which we decide what to keep and what to throw away in the raw IMU and muscle
signals.

\textbf{Where the signals live.} Deliberate hand and wrist motion has energy
mostly below about 5~Hz; physiological tremor sits near 8--12~Hz; mains hum is
at 50/60~Hz; and surface-EMG bursts from forearm muscles carry useful energy
roughly 20--450~Hz. These bands barely overlap, which is convenient: we can separate
them with the right magnitude responses. A low-pass on the accelerometer keeps
the gesture motion and rejects tremor and hum; a band-pass on the EMG keeps the
muscle bursts and removes both DC drift and high-frequency junk.

\textbf{Why phase matters here.} Gestures are short, sharp \emph{transients} ---
a flick, a tap, a clench. A filter with strongly nonlinear phase will smear that
transient in time, blurring exactly the onset we need for low-latency, musical
triggering. So we favor linear-phase (symmetric FIR) filters near the signal
path, accepting their fixed delay because at least that delay is the \emph{same}
for every frequency and therefore predictable.

\textbf{The hard limit.} A frequency response only \emph{re-weights} tones; it
cannot conjure information the sensor never captured. A wrist-mounted IMU alone
measures the gross motion of the whole forearm/wrist; no choice of $h[n]$ turns
it into a sensor of individual finger-joint angles --- that signal simply is not
present in the IMU stream. Detecting which finger moved needs forearm EMG or
force myography (FMG); reconstructing continuous finger pose needs sensors
\emph{on the fingers}. Filtering cleans up what the sensor saw; it cannot add a
sense the sensor does not have.
\end{goalbox}

\begin{buildbox}
On an Arduino or Pi reading an IMU at, say, $f_s = 200$~Hz, a cheap and
effective anti-tremor low-pass is a length-$M$ moving average plus a small
linear-phase trim, or a single-pole IIR smoother:
\begin{lstlisting}[language=C]
// One-pole low-pass: y[n] = a*x[n] + (1-a)*y[n-1]
// Cutoff (rad/sample) ~ a for small a; choose a from f_c, f_s:
//   a = 1 - exp(-2*pi*f_c/f_s)
float a = 1.0f - expf(-2.0f * 3.14159265f * 3.0f / 200.0f); // ~3 Hz cutoff
static float y = 0.0f;
float filt(float x) { y += a * (x - y); return y; }
\end{lstlisting}
This costs one multiply and one add per sample --- trivial on an 8-bit MCU.
Its frequency response is $H(e^{\jj\omega}) = a/(1 - (1-a)e^{-\jj\omega})$: a
gentle low-pass with a single pole at $1-a$. It is \emph{not} linear-phase
(its group delay rises near DC), so when transient timing is critical, prefer a
short symmetric FIR. Measure the actual response by feeding the firmware a
recorded chirp and running the \texttt{freq\_response} code above on the output.
\end{buildbox}

\section{Summary}
\label{ch14:summary}

\begin{itemize}
  \item Feeding an LTI system a complex exponential $e^{\jj\omega n}$ returns the
  same exponential scaled by a complex number $H(e^{\jj\omega})$. Complex
  exponentials are the eigenfunctions of LTI systems; $H$ collects the
  eigenvalues.
  \item $H(e^{\jj\omega})$ is the DTFT of the impulse response $h[n]$. Time and
  frequency descriptions hold identical information.
  \item $\abs{H(e^{\jj\omega})}$ is the gain at each frequency;
  $\angle H(e^{\jj\omega})$ is the phase shift. A cosine in comes out a cosine
  of the same frequency, rescaled and phase-shifted.
  \item Group delay $\tau_g = -\dd(\angle H)/\dd\omega$ is the envelope delay.
  Linear phase $\Rightarrow$ constant group delay $\Rightarrow$ no waveform
  smearing.
  \item Filtering multiplies spectra: $Y = HX$. The filter ``paints'' the input
  spectrum with its gain curve; it cannot create new frequencies.
  \item A pole near the unit circle gives resonance: a tall narrow peak in
  frequency and a long decaying ring in time --- two faces of one system.
  \item Abrupt impulse responses give leaky spectra; smooth (windowed) ones give
  clean spectra. Symmetric taps give linear phase.
\end{itemize}

\begin{connectionbox}
\textbf{Where We Are and Where We Go Next}
\begin{itemize}
  \item \textbf{What we have so far}: $H(e^{\jj\omega})$ gives the frequency-domain face of any
  stable LTI system: complex exponentials pass through scaled by $H(e^{\jj\omega})$, convolution
  turns into pointwise multiplication $Y = HX$, resonances show up as tall peaks, and group delay
  $-\dd(\angle H)/\dd\omega$ measures envelope latency.
  \item \textbf{Where the unit circle stops}: Evaluating only on $\abs{z} = 1$ works when the
  Fourier sum converges, so it cannot directly handle growing transients, initial conditions, or
  unstable feedback loops whose poles have wandered outside the unit circle.
  \item \textbf{Next}: Chapter~15 steps off the unit circle onto the whole complex plane with the
  \textbf{$z$-Transform and the Laplace Transform} ($z = r e^{\jj\omega}$), turning stability,
  resonance, and notch placement into pole--zero geometry.
\end{itemize}
\end{connectionbox}

\section{Exercises}
\label{ch14:exercises}

\begin{exercise}
For the difference filter $h=[\tfrac12,-\tfrac12]$, the magnitude response is
$\abs{\sin(\omega/2)}$. (a) Find its value at $\omega=0$, $\pi/2$, and $\pi$.
(b) At what $\omega$ has the gain reached $0.707$ (the $-3$~dB point)?
(c) Send $x[n]=\cos(\pi n)$ (the fastest sequence) through it and verify the
output amplitude matches $\abs{H(e^{\jj\pi})}$.
\end{exercise}

\begin{exercise}
A pure delay of $d$ samples has impulse response $h[n]=\delta[n-d]$. Compute
$H(e^{\jj\omega})$. Show that its magnitude is $1$ at every frequency and its
group delay is exactly $d$ at every frequency. Explain in words why a pure delay
is the ideal linear-phase ``filter'' that changes nothing but timing.
\end{exercise}

\begin{exercise}
A three-tap symmetric filter is $h=[\tfrac14,\tfrac12,\tfrac14]$. (a) Show that
$H(e^{\jj\omega}) = e^{-\jj\omega}\big(\tfrac12 + \tfrac12\cos\omega\big)$.
(b) Read off the magnitude and phase. (c) What is the group delay? (d) Is this
low-pass or high-pass, and what is its gain at DC and at Nyquist?
\end{exercise}

\begin{exercise}
Two filters have the same magnitude response but different phase: filter A is
$h_A=[1,0,0]$ delayed appropriately, and filter B reverses the taps of some
asymmetric $h$. Construct a concrete length-3 example where $\abs{H_A}=\abs{H_B}$
for all $\omega$ but the filters delay a transient differently. (Hint: compare
$h=[1, 0.5, 0]$ with its time-reversal.)
\end{exercise}

\begin{exercise}
You sample a sensor at $f_s = 500$~Hz and want to remove a $50$~Hz hum.
(a) What digital frequency $\omega_0$ (in radians/sample) does $50$~Hz map to?
(b) Sketch the magnitude response you would want (flat, with a notch at
$\pm\omega_0$). (c) Why must the notch be placed at both $\pm\omega_0$ rather than
only $+\omega_0$ for a real-valued filter?
\end{exercise}

\begin{exercise}
For the single-pole low-pass $y[n]=a\,x[n]+(1-a)\,y[n-1]$ with $0<a<1$:
(a) Find $h[n]$ (the impulse response) in closed form.
(b) Find $H(e^{\jj\omega})$ and show $\abs{H(e^{\jj0})}=1$.
(c) Explain how the cutoff moves as $a\to 0$ versus $a\to 1$, and connect this
to the length of the impulse response.
\end{exercise}

\begin{exercise}
A resonator has poles at radius $r$ and angle $\omega_0$. (a) As $r\to 1$,
describe what happens to the height and width of the spectral peak and to the
decay time of the impulse response. (b) The half-power bandwidth of the peak is
roughly $\Delta\omega \approx 2(1-r)$ for $r$ near $1$. If $r=0.99$, estimate the
bandwidth and the number of samples the ring takes to decay to $1/e$. Comment on
how ``sharp in frequency'' and ``long in time'' are linked.
\end{exercise}

\begin{exercise}
Using the \texttt{freq\_response} code in the chapter, plot the magnitude (in dB)
of the length-5, length-11, and length-21 box-average filters on the same axes.
(a) How does the cutoff frequency change with length? (b) Does the first
side-lobe level (about $-13$~dB) change with length? (c) Explain why making the
filter longer sharpens the transition but does not fix the leaky side lobes ---
and name the chapter that does fix them.
\end{exercise}

\begin{exercise}
A wrist IMU streams accelerometer data at $200$~Hz. Deliberate gesture motion is
below $\sim 5$~Hz; tremor is near $\sim 10$~Hz. (a) Choose a low-pass cutoff in
hertz that keeps gestures but suppresses tremor, and convert it to $\omega_c$.
(b) The one-pole smoother in the build box has rising group delay near DC.
Explain, in terms of musical triggering latency, why this could be a problem and
when a symmetric FIR would be the better choice. (c) State plainly one thing this
filtered IMU signal still cannot tell you about the hand, and what sensor you
would add to get it.
\end{exercise}

\begin{exercise}
(Conceptual.) Explain, without equations, why an LTI filter can never turn a
single input sine into two output sines of different frequencies. Then give a
real-world example of a \emph{non}-linear device that does create new frequencies,
and say how you would recognize its signature in a spectrum.
\end{exercise}
